
\subsection{实验结果与分析}

\begin{table}[ht]
\centering
\caption{开发集对照（每组 37 题）}
\label{tab:fqa-train}
\begin{tabular}{@{}lrrrr@{}}
\toprule
方案 & 支持 & 矛盾 & 信息不足 & 支持率\\
\midrule
基础 & 30 & 6 & 1 & 81.1\%\\
增强 & 33 & 4 & 0 & 89.2\%\\
\bottomrule
\end{tabular}
\end{table}

\textbf{开发集。} 两种方案均答完 37 题。增强方案多获 3 个支持判定，支持率提高 8.1 个百分点；模型总用量从 513\,167 降到 461\,295 token，减少约 10.1\%。模型调用由 141 次降至 137 次，工具调用由 246 次降至 234 次。在相同并发数下，批次用时分别为 30.02 和 26.70 秒；平均单题用时为 5.57 和 5.28 秒。耗时包含网络等待，是本次运行的观测值。索引从 9 份文件得到 4\,288 个片段，构建耗时 2.32 秒，索引约 12.85 MiB，两组共同使用。

初版增强方案有一题因模型输出未转义换行而未答。修复解析后，两种配置重跑开发集，预测文件仍是标准 JSON。基础方案有一次评测超时，按原答案、提示和模型补判。增强方案多获 3 个支持判定，入选正式测试。

\textbf{可核查的答案。} Snowflake 营收题中，模型查年度记录、计算增速，再检索管理层描述。日志中的 \texttt{F2509703} 和 \texttt{F2509712} 分别给出截至 2024 和 2025 年 1 月末的财年营收：28.065 亿美元和 36.264 亿美元。计算得到同比增长 29.21\%，低于上一年的 35.86\%。因此答案把“营收继续增加”和“增速放缓”分开叙述。

\textbf{失败样例。} Moderna 业务构成题中，增强方案引用了年报中“一个可报告分部”的表述，却没有说清哪类产品贡献最多收入。参考答案按产品业务线划分。遇到这类问题，应同时说明会计分部和产品线。


\begin{table}[ht]
\centering
\caption{正式测试成绩（增强方案）}
\label{tab:fqa-test}
\begin{tabular}{@{}lrrrr@{}}
\toprule
公司 & 题数 & 支持 & 矛盾 & 支持率\\
\midrule
Palantir & 10 & 10 & 0 & 100.0\%\\
Abbott   & 9 & 7 & 2 & 77.8\%\\
合计     & 19 & 17 & 2 & 89.5\%\\
\bottomrule
\end{tabular}
\end{table}

\textbf{正式测试。} 冻结的增强方案自动完成 19 题，格式检查确认没有缺失或空白答案。Pro 对全部答案独立评分，无信息不足或接口失败；按公司等权平均，支持率为 88.9\%。作答共调用模型 73 次和工具 125 次，4 次工具请求返回错误，均未导致空答。Flash 消耗 238\,132 token，含输入 226\,958 和输出 11\,174 token；批次用时 15.24 秒，单题平均 5.47 秒。

两道矛盾题都涉及 Abbott 的分部。其一，答案用诊断业务的分部总额衡量趋势，检查器则依据参考答案判定诊断业务也是近期增长来源。回查 2024 年报附注 16（日志 \texttt{H3493}），诊断分部净销售额由 99.88 亿美元降至 93.41 亿美元，说明“总额净下降”时还应与内部产品线和一次性需求的走势分开。该答案另把合并净销售额增幅写成 2.5\%；由同表 419.50 亿和 401.09 亿美元重算，应为约 4.6\%。编号真实并不能防止模型抄错比例。

其二，医疗器械占分部销售额约 45.3\%，其余三块合计约 54.7\%。答案把“最大板块”直接概括成“不均衡”，与参考答案强调的多元化发生冲突。报告记录这些差异用于分析；首次正式预测和原评分均保留，未据测试反馈改动方法。


\subsection{结论与局限}
增强配置在开发集以较少调用取得更高支持率，但改进成组进行，无法量化各组件的贡献。测试只覆盖两家公司；检查器的支持判定也不能替代对金额、期间和口径的人工核验。

结果附件含冻结配置、预测、真实调用日志和课程评测原件。仓库运行说明提供复现命令。
